\documentclass[answers,12pt,addpoints]{exam} 
\usepackage{import}

\import{C:/Users/prana/OneDrive/Desktop/MathNotes}{style.tex}

% Header
\newcommand{\name}{Pranav Tikkawar}
\newcommand{\course}{Spatial Statistics}
\newcommand{\assignment}{Homework 2.1}
\author{\name}
\title{\course \ - \assignment}

\begin{document}
\maketitle

\begin{questions}
    \question I have been looking more into the Matern covariance and I was wondering what was the origin/history of this function. I was curious to know if it was created/chosen for a specific application or if it was just found to be a useful class of covariance functions after the fact. Honestly, i feel like it has a super convoluted form and thus should probably be originated for a specific reason either mathematically or application wise.

    \question I find the semivariogram to be an interesting extension of the covariance function for spatial data. I realized that the notation of having evenly spaced data points is not always realistic in spatial statistics. I was curious to know if there are any methods of transforming irregularly spaced data into evenly spaced data via some "changing of basis" with a jacobian that is commonly used in spatial statistics.

    \question The notion of Fixed-domain asymptotic is interesting to me. I was hoping you could elaborate more on this. While the definition and apparent application is obvious, I was curious to know if there are a similar class of asymptotic theorems for these fixed-domain scenarios (with the exception of Stationarity). Especially since many of the classical asymptotic results in statistics rely on iid assumptions which are clearly violated in spatial statistics. 

    \question This method of making a predictor using the graph is actually a better start than I thought. When I was working on my research, I was hell bent on trying to model the best predictor for some really convoluted model, but by graphing behavior across a large number of simulations, I should get a better idea of what the predictor should look like.

    \question I think the part of pg 34 which highlights the strict requirements for a staitionary, gaussian process when dealing with spacial data clears up some of my confusion around what type of functions can be used to model spatial data, but not all. While I still want to explore more general $L_2$ processes more thoroughly, I see how the gaussian process is sufficient and simple to model complex spatial data. (all spatial data is complex)

    \question Is there analog to the Nyquist-Shannon sampling theorem for spatial data? Honestly, I think it might just be the same thing, but I was curious if there was a different name for it in spatial statistics.

    \question This was a significantly more confusing chapter for me. I will spend more time going over it. I was hoping to get a better understanding of periodograms and semivariograms. While I understand the definitions and how to compute them, I was hoping to get a better understanding of how they are used in practice other than as exploratory data analysis tools. 
    \end{questions}

\end{document}